\subsection{Large molecular datasets with released partitions}
\label{sec:native-molecular}

Released benchmark partitions provide a large-data comparison beyond the
ChEMBL scaling study. PCQM4Mv2, MOSES and ogbg-molpcba supply molecular SMILES
and native assignments \citep{hu2021ogblsc,polykovskiy2020moses,hu2020ogb}.
Retrieving only the required metadata used \NativeMetadataMB{} MB, including
archive byte ranges. No coordinate files or processed graph bundles were downloaded.
The datasets contain 3.75 million, 1.94 million and 438,000 original records.
All methods use the same successfully parsed rows; the supplement records
the 39 PCQM4Mv2 and two ogbg-molpcba parsing failures and retained identities.

\begin{table}[htbp]\centering\small
\begin{tabular}{lrrrrrr}
\toprule
Dataset & $n$ & Assignment & Global & Nearest & Seconds & GiB\\\midrule
PCQM4Mv2 & 3,746,581 & Native & 0.978 & 0.710 & -- & --\\
PCQM4Mv2 & 3,746,581 & Exact assignment & 0.620 & 0.537 & 304.0 & 55.4\\
PCQM4Mv2 & 3,746,581 & 1\% corridor & 0.619 & 0.536 & 300.5 & 55.4\\
MOSES & 1,936,962 & Native & 1.000 & 0.723 & -- & --\\
MOSES & 1,936,962 & Exact assignment & 0.784 & 0.605 & 186.0 & 28.9\\
MOSES & 1,936,962 & 1\% corridor & 0.783 & 0.598 & 185.3 & 28.9\\
ogbg-molpcba & 437,927 & Native & 0.963 & 0.651 & -- & --\\
ogbg-molpcba & 437,927 & Exact assignment & 0.781 & 0.570 & 40.6 & 7.0\\
ogbg-molpcba & 437,927 & 1\% corridor & 0.779 & 0.564 & 40.3 & 7.0\\
\bottomrule\end{tabular}
\caption{Released partitions and low-rank controls on the same parsed records. Global is overlap relative to count-matched random. Nearest weights role means from 128 sampled queries against the complete training fold. Exact assignment preserves native counts and omits local refinement; the corridor uses 1\% tolerance and the standard size cutoff. Low-rank rows average three seeds. Seconds cover the public call; GiB is peak process RSS before auditing. Native construction costs are unavailable. Parsing losses, per-role scores and sampling intervals are in the supplement.}
\label{tab:native-molecular}
\end{table}


On \NativeLargestRecords{} retained PCQM4Mv2 records, exact-count assignment
optimization changes mean global overlap from \NativePcqmNativeGlobal{} to
\NativePcqmExactGlobal{} and sampled nearest-training exposure from
\NativePcqmNativeNearest{} to \NativePcqmExactNearest{}. It takes
\NativePcqmExactSeconds{} seconds on 16 CPU cores with mean peak process RSS
\NativePcqmExactMemory{} GiB. This assignment-only control explicitly omits
refinement. The corridor setting uses the standard size-dependent refinement
rule; point refinement is skipped above two million records. These are
separate protocols from the two-million-entity scaling measurements above.

The released partitions serve different purposes. PCQM4Mv2 uses CID-based
roles, MOSES has distinct random and scaffold tests, and ogbg-molpcba
uses scaffolds. These roles remain separate in the evaluation. The low-rank
controls optimize aggregate overlap with native counts or a declared corridor;
they do not preserve every native chemical grouping constraint. Their lower
similarities therefore describe alternative structural boundaries rather than
an isolated comparison of optimizers on identical native tasks. Original
native construction times were not supplied, and loading times are recorded
separately. Per-role exposure, Monte Carlo intervals and sampled returned-cut
fidelity accompany the full outcome inventory in the supplement.
